% page 447 du fac-simile (image p-446 du scan)
% bloc 162.6 x 246.3 mm ; marge gauche 39.4 mm
\pgc{17.780mm}{0.000mm}{
\l{\cel{52}439}
\l{}
\l{\sou{pilono}. (\sou{arkitekt.}) I. (en\cel{1}\sou{Egiptia, antique}).Sorto di avan-}
\l{\cel{3}korpo qua havas la formo di piramido quadriangula trunka, e}
\l{\cel{3}kun porda aperturo centrala, quan la Egiptiani antiqua kustu-}
\l{\cel{3}mis erektar avan la enireyo di sua templi. - II. Dekor-motivo}
\l{\cel{3}qua havas la formo di pilastro quadriangula, lokizita an la}
\l{\cel{3}lateri di enireyo. - DEFIRS.}
\l{}
\l{\sou{piloro}. (\sou{anat}.) Orifico infra di la stomako, tra qua la nutrivi}
\l{\cel{3}pasas aden la intestini. - EFIS.}
\l{}
\l{\sou{pilorio}. Fosto an qua onu ligis per fero-ringo la persono qua}
\l{\cel{3}esis kondamnita ad expozeso publika (kom puniso legala). -}
\l{}
\l{}
\l{\sou{pilotar}. (\sou{trans.}) Duktar navo en la mar-paseyi desfacile tra-}
\l{\cel{3}irebla, ye la enireyo od ekireyo di portuo. - DEFIS.}
\l{}
\l{\sou{pilulo}. (\sou{farmacio)}.\cel{2}Medikamento fasonita a mikra globo, maxim-}
\l{\cel{3}multa-kaze envelopita ye folio dinega ek oro od arjento, qua}
\l{\cel{3}impedas lua adhero e plufaciligas lua engluteso. - DEFIRS.}
\l{}
\l{\sou{pimento}. I.Drinkajo ek mielo, spici e vino. - II. (\sou{bot.}) Bero}
\l{\cel{3}konoforma, kun surfaco glata e briloza ; komence bele-verda,}
\l{\cel{3}lu divenas (lor lua matureso) brile-reda - uzata kom kondi-}
\l{\cel{3}mento. - L.\cel{1}\sou{mirtus pimenta}. - DEFIS.}
\l{}
\l{\sou{pimpinelo}. (\sou{bot.}) Planto herbatra, aromata, ek la familio \textquotedbl{}roza-}
\l{\cel{3}cei\textquotedbl{}, uzata kom kondimento. - L.\cel{1}\sou{pimpinella}. -}
\l{}
\l{\sou{pino}.\cel{2}(\sou{bot.)}\cel{2}Granda arboro rezinoza, sempre verda, ek la fa-}
\l{\cel{3}milio \textquotedbl{}koniferi\textquotedbl{}. - L.\cel{1}\sou{pinus sylvestris}. - EFIS.}
\l{}
\l{\sou{pinaklo}. (\sou{arkitekt.})\cel{2}Tekto-karpenturo dekorita, e finanta per}
\l{\cel{3}pinto. - DEFIS.}
\l{}
\l{\sou{pinastro}. (\sou{bot.})\cel{2}Pino sovaja (mar-litorala) de qua onu extrak-}
\l{\cel{3}tas terpentino. - L.\cel{1}\sou{pinus pinaster}. - EFIS.}
\l{}
\l{\sou{pinco}.\cel{2}(\sou{tekn.}) Mikra te oi\cel{2}o,uzata diverse en la mestieri, en}
\l{\cel{3}kirurgio. - \textquotedbl{}Pinco\textquotedbl{}diferas de \textquotedbl{}tenalio\textquotedbl{} per to ke lu esas}
\l{\cel{3}utensilo sen charniro (la du parti ye qua lu konsistas esas}
\l{\cel{3}asemblita per\cel{2}resorto o parto flexebla). - DEFIRS.}
\l{}
\l{\sou{pinchar}. (\sou{trans.}) Klemar forte inter la fingri. - EF.}
\l{}
\l{\sou{pinealo}. (\sou{anat.)}\cel{2}Mikra korpo, avan la cerebelo. - DEFIS.}
\l{}
\l{\sou{pinglo}. Mikra stango de latuno, pinta ye un extremajo, kapoza ye}
\l{\cel{3}la altra, quan onu uzas por ligar. - EF.}
\l{}
\l{\sou{pinguino}. (\sou{zool.)}\cel{3}Ucelo palmipeda di la Nordo-mari, kun ali}
\l{\cel{3}kurta. - L.\cel{1}\sou{aptenodytes}. - DEFIRS.}
\l{}
\l{pinio. (bot.) Sorto di pino qua kreskas en la regioni sudala}
\l{\cel{3}di Europa, kun kapo parasuno-forma. - L.\cel{1}\sou{pinus pinea}.- DF.}
}
